% ===========================================================================
% Resume bullets -- RL Cloud Scheduler project
%
% Copy ONE version into your resume. Percent signs and underscores are escaped.
% Version A: portable (standard itemize, works in any template)
% Version B: Jake's Resume macros (\resumeProjectHeading, \resumeItem)
% Version C: compact 3-bullet version for a one-page resume
%
% Every number below is traceable to results/summary_table.csv,
% results/full_comparison.csv, results/ga_logs/ and tests/.
% ===========================================================================


% ---------------------------------------------------------------------------
% Version A -- portable
% ---------------------------------------------------------------------------
\section*{Projects}
\textbf{RL Cloud Scheduler} \hfill 2026\\
\textit{Python, PyTorch, Stable-Baselines3, Gymnasium, React, pytest}
\begin{itemize}
  \item \textbf{Simulator + benchmark harness:} built an event-driven
        cloud-scheduling environment (10 servers, 2{,}000-job scenarios derived
        from the Google ClusterData2011-2 trace) and a shared
        \texttt{policy(obs, env)} interface; benchmarked 9 schedulers across
        heuristic, soft-computing, and reinforcement-learning families
        (27 policy--scenario settings, 48 seeded runs).

  \item \textbf{Fuzzy scheduler + genetic tuning:} designed a zero-order Sugeno
        fuzzy inference system (3 linguistic inputs, 18 rules) and a 24-gene
        genetic algorithm that tunes membership breakpoints and rule
        consequents; enforced training-only model selection with an
        assertion-based data-leakage guard and 3 independent GA seeds.

  \item \textbf{RL implementation:} implemented DQN, Double DQN, PPO, and A2C
        with Stable-Baselines3 and evaluated every policy through one harness
        on average/p95 wait, CPU utilization, completion rate, invalid actions,
        and makespan.

  \item \textbf{Result:} showed RL was not necessary for this workload -- the
        fuzzy scheduler matched the best heuristic within 0.53\% and posted the
        lowest bursty-load wait (46{,}869\,s), while the strongest RL result was
        still 21.8\% slower; traced the gap to invalid placements (up to
        $\sim$2{,}000 per episode for DQN/PPO).

  \item \textbf{Engineering:} built a React replay dashboard with a step-by-step
        fuzzy decision inspector and 6 report figures, and added a 33-test
        pytest suite covering the environment, Double DQN, fuzzy inference, and
        GA operators; kept the pipeline reproducible from raw trace to dashboard.
\end{itemize}


% ---------------------------------------------------------------------------
% Version B -- Jake's Resume (https://github.com/jakegut/resume)
% Use only if your template defines these macros.
% ---------------------------------------------------------------------------
\resumeProjectHeading
  {\textbf{RL Cloud Scheduler} $|$ \emph{Python, PyTorch, Stable-Baselines3, Gymnasium, React}}{2026}
  \resumeItemListStart
    \resumeItem{Built an event-driven cloud-scheduling simulator (10 servers,
      2{,}000-job scenarios from Google ClusterData2011-2) and a shared
      \texttt{policy(obs, env)} interface; benchmarked 9 schedulers across
      heuristic, soft-computing, and RL families over 27 policy--scenario
      settings (48 seeded runs).}
    \resumeItem{Designed a zero-order Sugeno fuzzy scheduler (3 inputs, 18
      rules) and a 24-gene genetic algorithm tuning membership breakpoints and
      consequents, with training-only selection, an assertion-based leakage
      guard, and 3 independent seeds.}
    \resumeItem{Implemented DQN, Double DQN, PPO, and A2C and evaluated all
      policies on average/p95 wait, CPU utilization, completion rate, invalid
      actions, and makespan through one shared harness.}
    \resumeItem{Showed RL was unnecessary: the fuzzy scheduler matched the best
      heuristic within 0.53\% and had the lowest bursty-load wait (46{,}869\,s);
      the strongest RL run was 21.8\% slower, with the gap traced to invalid
      placements (up to $\sim$2{,}000 per episode for DQN/PPO).}
    \resumeItem{Built a React replay dashboard with a fuzzy decision inspector
      and 6 report figures; added a 33-test pytest suite and kept the pipeline
      reproducible end to end.}
  \resumeItemListEnd


% ---------------------------------------------------------------------------
% Version C -- compact (3 bullets)
% ---------------------------------------------------------------------------
\begin{itemize}
  \item \textbf{RL Cloud Scheduler} --- built an event-driven 10-server
        simulator from Google ClusterData2011-2 and a unified
        \texttt{policy(obs, env)} harness; benchmarked 9 policies (3 heuristics,
        FIS/GA-FIS, DQN/Double DQN/PPO/A2C) over 27 policy--scenario settings
        and 48 seeded runs.
  \item Designed an 18-rule Sugeno fuzzy scheduler plus a 24-gene genetic
        algorithm for membership/rule tuning, with training-only selection and
        an assertion-based leakage guard.
  \item Demonstrated that the fuzzy scheduler matched the best heuristic within
        0.53\% and beat every RL agent (best RL run 21.8\% slower), then shipped
        a React decision-inspector dashboard and a 33-test suite.
\end{itemize}


% ---------------------------------------------------------------------------
% Optional skills line
% ---------------------------------------------------------------------------
\textbf{Skills:} Python, NumPy, pandas, PyTorch, Stable-Baselines3, Gymnasium,
fuzzy logic, genetic algorithms, React, Matplotlib, pytest, Git.
